\Article{26}
When the intensity of terrestrial magnetism $T$ is expressed by a \textit{number} $k$, that number is based on a particular unit $V$, namely a force of the same kind. Its relation to the other units given directly is contained in what precedes, though in a somewhat complicated form. It is therefore worth setting out this relation again here, so that we can see in elementary terms how the number $k$ changes if we start from different fundamental units.

\SourceLeaf{120}
\expandafter\def\csname HistoricalNoteText13\endcsname{It will hardly be necessary to point out that the meanings previously assigned to the letters no longer apply here.}
To establish the unit $V$, we had to start from the unit of free magnetism $M$\HistoricalNote{13} and the unit of distance $R$. We set $V$ equal to the force exerted by $M$ at distance $R$.

As the unit $M$ we took the quantity of magnetic fluid that, acting on an equal quantity $M$ at distance $R$, produces a motive force (or a pressure, if you prefer) equal to $W$, the force taken as the unit. That is, it equals the force exerted by the accelerating force $A$ taken as the unit on the mass $P$ taken as the unit.

There are two ways to establish the unit $A$. It can be taken from a similar force given directly, for example from gravity at the place of observation, or from its effect in moving bodies. The latter method, which we followed in our calculations, requires two new units, the unit of time $S$ and the unit of velocity $C$. We take as the unit $A$ the accelerating force that, acting for a time $S$, produces velocity $C$. Finally, we take this velocity to be that of uniform motion over a distance $R$ in time $S$.

Thus the unit $V$ clearly depends on three units, either $R,P,A$ or $R,P,S$.

Suppose we now replace the units $V,R,M,W,A,P,C,S$ with $V',\allowbreak R',\allowbreak M',\allowbreak W',\allowbreak A',\allowbreak P'.\allowbreak\ C',\allowbreak S'$, related to each other in the same way as the previous ones. Suppose also that in the unit $V'$ terrestrial magnetism is expressed by the number $k'$. We must find how it is related to $k$.

If we set
\[
\begin{aligned}
V&=vV'\\
R&=rR'\\
M&=mM'\\
W&=wW'\\
A&=aA'\\
P&=pP'\\
C&=cC'\\
S&=sS'
\end{aligned}
\]
then $v,r,m,w,a,p,c,s$ are pure numbers, and we have
\SourceLeaf{121}
\[
\begin{aligned}
kV&=k'V'\quad\text{or}\quad kv=k'\\
v&=\frac m{rr}\\
\frac{mm}{rr}&=w=pa\\
a&=\frac cs\\
c&=\frac rs
\end{aligned}
\]
Combining these equations gives
\[
\begin{aligned}
\text{I.}\qquad k'&=k\BareRoot{\frac p{rss}}\\
\text{II.}\qquad k'&=k\BareRoot{\frac{pa}{rr}}
\end{aligned}
\]
As long as we keep to the method we followed in our calculations, we must use the first formula. If, for example, we take the meter and gram as units in place of the millimeter and milligram, then $r=\frac1{1000}$ and $p=\frac1{1000}$, so that $k'=k$. If we take the Paris line and the Berlin pound, we have $r=\frac1{2{,}255829}$ and $p=\frac1{467711{,}4}$, so that $k'=0{,}002196161k$. Thus, for example, experiments VIII yield the value $T=0{,}0039131$.

If we prefer to follow the other method and take gravity as the unit of accelerating forces, we set $a=\frac1{9811{,}63}$ for the Göttingen observatory. As long as we keep the millimeter and milligram, the numbers $k$ must then be multiplied by $0{,}01009554$, and any changes to those units must be handled according to formula II.

